\section{Objetivo general}
Identificar si el rechazo hacia la IA generativa viene del conocimiento del uso de inteligencia artificial o de la calidad del contenido, mediante la creación y comparación de dos versiones de un mismo videojuego para evaluar la percepción de los estudiantes universitarios.

\section{Objetivos específicos}
\begin{itemize}
\item Garantizar la estabilidad de la prueba, desarrollando e implementando ambas versiones sin errores críticos o que favorezcan una versión sobre la otra.
\item Comparar las diferencias en la experiencia de usuario (UX) reportadas por los estudiantes al interactuar de manera ciega con ambas versiones del juego.
\item Identificar las preferencias específicas de los jugadores respecto a componentes individuales como el diseño de niveles, las mecánicas de juego y el apartado técnico.    
\item Garantizar una diferencia funcional entre ambas versiones del videojuego, asegurando que se mantenga el concepto original para ambas versiones.
\item Analizar la correlación entre el conocimiento previo de los estudiantes sobre el uso de IA generativa en el desarrollo y su nivel de aceptación o rechazo hacia cada versión del videojuego.
\end{itemize}